The starting point for any Corvus calculation is to request a property or list of properties to calculate.
This is done using the keyword target\_list.
From here, corvus will look for available codes that can calculate each target, and figure out what input is needed.
There are many possible targets for Corvus, but here I will list a few of the most important ones for calculating spectra
using FEFF10, and how to use them. In general, when using FEFF10 an average over unique sites in the structure must be taken
in order to get the full spectrum. Thus in order to perform a spectral calculation using an external structure as input (the suggested method)
you must define,
\begin{verbatim}
target_list{ cfavg} # Calculate the average of some property.
cfavg_target{ xanes } # Average the XANES, could also be xes, rixs, exafs, ... or multiple properties.
\end{verbatim}

To specify the structure, you can use a variety of file formats, e.g., .xyz, .cif, XDATCAR, or you can have 
corvus download a structure from the Materials Project.
Example:
\begin{verbatim}
cif_input{ your_structure.cif }
\end{verbatim}

You will also need to specify which element and edge you want, e.g., for the Cu K-edge,
\begin{verbatim}
absorbing_atom_type{ Cu }
feff.edge{ K }
\end{verbatim}

Thus a minimal input file for the Cu K-edge of CuO would look as follows:
\begin{verbatim}
target_list{ cfavg }
cfavg_target{ xanes }
cif_input{ CuO.cif }
absorbing_atom_type{ Cu }
feff.edge{ K }
\end{verbatim}

There are of course many useful and important options available, but this will give you a calculation that has reasonable defaults.
One special target to note is the "loop" target which allows the user to defined loops over the values associate with any keyword or set 
of keywords. For single line keywords, the structure is relatively simple:
\begin{verbatim}
target_list{ loop } # Make corvus produce a loop over some parameters.
loop_target{ cfavg } # The target property to calculate for each calculation is an average

cfavg_target{ xanes } # Average the xanes over different unique sites in the structure.

# Now define the loop:
loop_parameter{
single_line_keyword
values_line1
values_line2
values_line3
...
}
# Optionally, define labels for the loop directories. Must have the same number of entries as loop_parameter.
# loop_labels{
# label1
# label2
# label3
# ...
# }

There are is also a special value for the first line after "loop_parameter" that loops over input files so that one can loop over 
any defined sets of keyword->value combinations, e.g., 
\begin{verbatim}
loop_parameter{
file
file1.in
file2.in
file3.in
...
}
\end{verbatim}
This will cause corvus to perform calculations with all parameters loaded from the main input file + any parameters existing in fileN.in. Input in fileN.in overrides input in the 
main input file.
